\documentclass[10pt]{article}
\usepackage{amsmath, fancyhdr, libertine}
\usepackage[libertine]{newtxmath}
\usepackage[margin=1in]{geometry}
\pagestyle{fancy}
\fancyhead[R]{\sc Mathematics Review}

\begin{document}
\section{Functions}
For each number $x$, a function $f$ assigns a unique number $y$ according to some rule.
\begin{equation*}
  \underbrace{y}_{\text{\shortstack{Dependent\\variable}}} = \underbrace{f}_{\text{Function}} \left( \underbrace{x}_{\text{\shortstack{Independent\\variable}}} \right)
\end{equation*}
For example, $y=f\left( x \right)=x^2$. The \emph{rule} to obtain $y$ is to take $x$ and square it. For example, $f\left( x=2 \right) = 4$, and $f\left( x=3 \right)=9$.

A variable can also be a function of more than one variable:
\begin{equation*}
  y=f\left( x_1,x_2 \right)=2x_1 + 3x_2
\end{equation*}
\section{Linear Functions}
A linear function is a function of the form:
\begin{equation*}
  y = a + bx
\end{equation*}
$a$ is the intercept on the vertical axis. $b$ is the slope. For example, $y=2+3x$.
\section{The Rate of Change}
\begin{itemize}
  \item Let $\Delta x = x^\prime - x$ be a \emph{small} change in $x$.
  \item A \emph{rate of change} is the ratio of two changes.
  \item If $y=f\left( x \right)$, the rate of change of $y$ with respect to $x$ is:
    \begin{equation*}
      \frac{\Delta y}{\Delta x} = \frac{f\left( x+\Delta x \right)- f\left( x \right)}{\Delta x}
    \end{equation*}
\end{itemize}
\subsubsection*{Examples:}
Linear function: $y=a+bx$.
  \begin{equation*}
    \frac{\Delta y}{\Delta x} = \frac{f\left( x+\Delta x \right)- f\left( x \right)}{\Delta x}=\frac{\left[ a + b\left( x+\Delta x \right)-\left[ a-bx \right] \right]}{\Delta x} = \frac{b\Delta x }{\Delta x} = b
  \end{equation*}
Quadratic function: $y=x^2$.
  \begin{equation*}
  \frac{\Delta y}{\Delta x} = \frac{f\left( x+\Delta x \right)- f\left( x \right)}{\Delta x}=\frac{\left( x+\Delta x \right)^2 - x^2}{\Delta x}=\frac{x^2 + 2x\Delta x+\left( \Delta x \right)^2 - x^2}{\Delta x} = \frac{2x\Delta x+\left( \Delta x \right)^2}{\Delta x}=2x + \Delta x
  \end{equation*}
If $\Delta x$ is very small, this is just $2x$.
\section{Derivatives}
The derivative of a function is defined as:
  \begin{equation*}
    \frac{df\left( x \right)}{dx} = \underset{\Delta x \rightarrow 0}{\lim}\frac{f\left( x+\Delta x \right)- f\left( x \right)}{\Delta x}
  \end{equation*}
  It is the rate of change as the change goes to zero. We write it as $\frac{dy}{dx}$ or $f^\prime\left( x \right)$.
\subsubsection*{Examples:}
\begin{align*}
  f\left( x \right)&= a x^b
  & \Longrightarrow\qquad
  \frac{df\left( x \right)}{dx} &= abx^{b-1} \\
  f\left( x \right)&= bx
  & \Longrightarrow\qquad
  \frac{df\left( x \right)}{dx} &= b \\
  f\left( x \right)&= a
  & \Longrightarrow\qquad
  \frac{df\left( x \right)}{dx} &= 0 \\
  f\left( x \right)&= \frac{1}{x} =x^{-1}
  & \Longrightarrow\qquad
  \frac{df\left( x \right)}{dx} &= -x^{-2} = -\frac{1}{x^2}
\end{align*}
The derivative evaluated at a particular point is the slope of the tangent at that point.
For linear functions, the derivative is the same everywhere and the slope of the tangent is always the same. For nonlinear functions, the slope of the tangent changes.
\section{Second Derivatives}
The second derivative of a function is simply the derivative of the derivative. We write it as $\frac{d^2y}{dx^2}$ or $f^{\prime\prime}\left( x \right)$.
\subsubsection*{Example:}
\[
  f\left( x \right) = x^3   \qquad \Longrightarrow \qquad  \frac{df\left( x \right)}{dx} = 3x^2  \qquad \Longrightarrow \qquad  \frac{d^2f\left( x \right)}{dx^2} = 6x
\]
\begin{itemize}
  \item The 2nd derivative measures the curvature of a function.
  \item A negative 2nd derivative means a function is concave.
  \item A positive 2nd derivative means a function is convex.
\end{itemize}

\section{Optimization}
Suppose $y=f\left( x \right)$ and we wanted to know what value of $x$ maximizes $f\left( x \right)$.

\begin{itemize}
  \item An extreme point of a function occurs when its derivative equals zero,
    i.e.~the extreme point of $f\left( x \right)$ is at the $x^\star$ that
    solves
    $f^\prime\left( x^\star \right)=0$.
  \item An extreme point can be either a maximum or a minimum (\emph{Draw examples}).
  \item A function $f\left( x \right)$ is maximized at $x^\star$ if
    $f^\prime\left( x^\star \right)=0$ and $f^{\prime\prime}\left( x \right)
    < 0$.
  \item A function $f\left( x \right)$ is minimized at $x^\star$ if
    $f^\prime\left( x^\star \right)=0$ and $f^{\prime\prime}\left( x \right)
    > 0$.
\end{itemize}
\emph{Example:} What value of $x$ maximizes $f\left( x \right)=2 + 4x - x^2$?
The first derivative is:
\[
  f^\prime\left( x \right)=4 - 2x
\]
We need to find the value of $x^\star$ that solves $f^\prime\left( x^\star
\right)=0$. This is $4-2x^\star=0 \Longrightarrow x^\star=2$.

Check that it is a maximum:
\[
  f^{\prime\prime}\left( x^\star \right)= -2 < 0
\]
Since $f^{\prime\prime}\left( x^\star \right) < 0$, this is indeed a maximum.


\section{Partial Differentiation}
Partial differentiation is used when a function has more than one
variable/input. The partial derivative of $f\left( x_1,x_2 \right)$ with
respect to $x_1$ is the derivative of the function with respect to $x_1$,
\underline{holding $x_2$ fixed}.
\subsubsection*{Examples:}
\begin{align*}
  f\left( x_1,x_2 \right)&= x_1 x_2
  & \qquad\Longrightarrow\qquad
  \frac{\partial f\left( x_1,x_2 \right)}{\partial x_1} &= x_2 \\
  f\left( x_1,x_2 \right)&= x_1  + x_2
  & \qquad\Longrightarrow\qquad
  \frac{\partial f\left( x_1,x_2 \right)}{\partial x_1} &= 1 \\
  f\left( x_1,x_2 \right)&= a x_1^b x_2^c
  & \qquad\Longrightarrow\qquad
  \frac{\partial f\left( x_1,x_2 \right)}{\partial x_1} &= abx_1^{b-1}x_2^c \\
  f\left( x_1,x_2 \right)&= a x_1^b  + x_2^c
  & \qquad\Longrightarrow\qquad
  \frac{\partial f\left( x_1,x_2 \right)}{\partial x_1} &= abx_1^{b-1} \\
\end{align*}
\begin{itemize}
  \item Notice that we use $\partial$ instead of $d$ to emphasize that this is a partial derivative.
  \item Notice also that when a term is not multiplied by the variable with are differentiating over, it 
\end{itemize}




\end{document}
